\begin{problem}[2019年第36届全国中学生物理竞赛复赛][微孔竹板声学共振吸收]{107}
2016年9月，G20峰会在杭州隆重召开，其会议厅的装饰设计既展示出中国建筑的节能环保理念，又体现了浙江的竹文化特色。图a给出了其部分墙面采用的微孔竹板装饰的局部放大照片，该装饰同时又实现了对声波的共振吸收。竹板上有一系列不同面积、周期性排列的长方形微孔，声波进入微孔后导致微孔中的空气柱做简谐振动。单个微孔和竹板后的空气层，可简化成一个亥姆霍兹共振器，如图b所示。假设微孔深度均为 $l$、单个微孔后的空气腔体体积均为 $V_{0}$、微孔横截面积记为 $S$。声波在空气层中传播可视为绝热过程，声波传播速度 $v_{\mathrm{s}}$ 与空气密度 $\rho$ 及体积弹性模量 $\kappa$ 的关系为
\[
v_{\mathrm{s}} = \sqrt{\frac{\kappa}{\rho}}
\]
其中 $\kappa$ 是气体压强的增加量 $\Delta p$ 与其体积 $V$ 相对变化量之比
\[
\kappa = - \frac{\Delta p}{\Delta V / V} = - V \frac{\Delta p}{\Delta V}
\]
已知标准状态（273K，$1\mathrm{atm}=1.01\times10^{5}\mathrm{Pa}$）下空气（可视为理想气体）的摩尔质量 $M_{mol}=29.0\,\mathrm{g/mol}$，热容比 $\gamma=\frac{7}{5}$，气体普适常量 $R=8.31\,\mathrm{J/(K\cdot mol)}$。
\begin{subproblem}
求标准状态下空气的密度和声波在空气中的传播速度 $v_{s}$；
\end{subproblem}
\begin{subproblem}
求上述亥姆霍兹共振器吸收声波的频率（用 $v_{s}$、$S$、$l$、$V_{0}$ 表示）；
\end{subproblem}
\begin{subproblem}
为了吸收频率分别为 $120\,\mathrm{Hz}$ 和 $200\,\mathrm{Hz}$ 的声波，相应的两种微孔横截面积之比应为多少？
\end{subproblem}
\end{problem}